% ---------------------------------------------------------------------
% Anexo de Objetivos de Desarrollo Sostenible (ODS)
% ---------------------------------------------------------------------

\chapter{Relación del trabajo con los Objetivos de Desarrollo Sostenible}
\label{cap:anexo-ods}

De acuerdo con el artículo 10.3 de la Normativa de Trabajos de Fin de Grado y Trabajos de Fin de Máster de la UPV (Consejo de Gobierno, 21/07/2022) y su Anexo I, este anexo recoge el grado de relación del TFG con los Objetivos de Desarrollo Sostenible (ODS) de la agenda 2030.

La \autoref{tab:ods} recoge únicamente los ODS con una relación real y directa con el trabajo; el resto se marca como \enquote{No procede} en vez de forzar una relación tangencial.

\begin{table}[htbp]
\centering
\caption{Grado de relación del trabajo con los Objetivos de Desarrollo Sostenible.}
\label{tab:ods}
\begin{tabularx}{\textwidth}{Xp{0.09\textwidth}p{0.09\textwidth}p{0.09\textwidth}p{0.16\textwidth}}
\toprule
\textbf{Objetivo de Desarrollo Sostenible} & \textbf{Alto} & \textbf{Medio} & \textbf{Bajo} & \textbf{No\newline procede} \\
\midrule
ODS 1. Fin de la pobreza. & & & & X \\
ODS 2. Hambre cero. & & & & X \\
ODS 3. Salud y bienestar. & & & & X \\
ODS 4. Educación de calidad. & & & & X \\
ODS 5. Igualdad de género. & & & & X \\
ODS 6. Agua limpia y saneamiento. & & & & X \\
ODS 7. Energía asequible y no contaminante. & & & & X \\
ODS 8. Trabajo decente y crecimiento económico. & & & & X \\
ODS 9. Industria, innovación e infraestructuras. & X & & & \\
ODS 10. Reducción de las desigualdades. & & & & X \\
ODS 11. Ciudades y comunidades sostenibles. & & X & & \\
ODS 12. Producción y consumo responsables. & & & & X \\
ODS 13. Acción por el clima. & & & & X \\
ODS 14. Vida submarina. & & & & X \\
ODS 15. Vida de ecosistemas terrestres. & & & & X \\
ODS 16. Paz, justicia e instituciones sólidas. & & & & X \\
ODS 17. Alianzas para lograr objetivos. & & & & X \\
\bottomrule
\end{tabularx}
\end{table}

\section*{Descripción de la alineación con los ODS de mayor relación}

\begin{itemize}
\item \textbf{ODS 9 (Alto)}: el trabajo desarrolla e integra una solución tecnológica de navegación autónoma basada en visión para un vehículo aéreo no tripulado, combinando un sistema empotrado embarcado (computadora de borde, controlador de vuelo, sensores) con un pipeline de visión por computador (descriptores visuales profundos, emparejamiento geométrico) diseñado y validado empíricamente sobre hardware real. Este tipo de integración hardware-software es directamente extrapolable a aplicaciones industriales de inspección, monitorización de infraestructuras o logística mediante drones.
\item \textbf{ODS 11 (Medio)}: la navegación por visión sin depender de la señal GNSS es especialmente relevante en entornos donde esta señal se degrada o no está disponible, como los entornos urbanos densos (por reflexión y bloqueo de la señal entre edificios, el llamado \enquote{cañón urbano}) o durante emergencias en las que la infraestructura de posicionamiento puede fallar. Aunque este trabajo no se centra en un caso de uso urbano concreto, la tecnología desarrollada y validada es aplicable a la resiliencia de los sistemas de navegación aérea en ciudades y comunidades.
\end{itemize}

% ---------------------------------------------------------------------
